\honest{Wrong Questions}{Eliezer Yudkowsky}{2008}

Where the mind runs against the grain of reality, it produces wrong questions. They cannot be answered on their own terms, only dissolved, by understanding the mental process that makes them feel like questions. One good cue: you cannot imagine any concrete state of the world that would answer the question. \nb{Wittgenstein's \textsc{Tractatus} (1921) says something close: ``If a question can be put at all, then it can also be answered.''}

Take the tree that falls unheard. No state of affairs makes ``sound'' really mean vibrations rather than experiences, and if you say ``it is the state of affairs where `sound' means acoustic vibrations,'' taboo the word ``means'' and try again.

If that is too easy, take free will: ``What concrete state of affairs, whether in deterministic physics, or in physics with a dice-rolling random component, could ever correspond to having free will?'' \nb{Hume named one in 1748: ``a power of acting or not acting, according to the determinations of the will,'' which belongs to anyone ``not a prisoner and in chains.'' Deterministic physics allows it. In ``Thou Art Physics,'' three months after this post, Yudkowsky gave a compatibilist answer of the same kind.}

If that is too easy, ask why anything exists at all. I do not know the answer. But from ``my previous experience with unanswerable questions,'' which I do not describe, I predict that the answer will be no grand First Cause but an insight showing that the question was wrong.

Mystery exists in the mind, not in reality, and ``Confusion exists in the map, not in the territory.'' Unanswerable questions mark places where the mind runs skew to reality. Magic does not enter the universe there. Such questions must be dissolved. Trying to answer them ``inevitably'' produces the worst kind of mysterious answer: seemingly strong arguments, but no new predictions even in hindsight, and the phenomenon as sacredly inexplicable as before. My one example is an answer I make up myself: perhaps nothing exists. Believing it would leave you as confused as before.

The wonderful thing about unanswerable questions is that ``they are always solvable, at least in my experience.'' \nb{Of the three examples, the falling tree had been treated in earlier posts; free will was set as homework the same day; for existence the author does not know the answer.} What Queen Elizabeth I thought on waking on her fortieth birthday is different: I can easily imagine answers, but the information is lost. Why anything exists seems so impossible that I infer I am confused, and that the truth is probably not very complicated, and will be visible once the confusion goes. ``This may seem counterintuitive if you've never solved an unanswerable question, but I assure you that it is how these things work.'' Tomorrow I give a simple trick for wrong questions.
